%% =====================================================================
%%  APENDICE B - LISTA DE VERIFICACAO DE ENCERRAMENTO
%% =====================================================================

\chapter{Lista de verificação de encerramento aplicada}
\label{ap:checklist}

Este apêndice aplica ao estudo de caso a lista de verificação de encerramento do
framework, com catorze itens. Para cada item registra-se a situação observada e a
evidência correspondente. A coluna ``Aplicável'' distingue os itens que não se
aplicam ao tipo de artefato produzido --- distinção necessária porque uma lista de
verificação com itens permanentemente inaplicáveis deixa de ser um instrumento de
controle.

{\footnotesize
	\begin{longtable}{@{}c p{6.2cm} c p{3.6cm}@{}}
		\caption{Lista de verificação de encerramento aplicada ao estudo de caso}
		\label{tab:checklist} \\
		\toprule
		\textbf{\#} & \textbf{Item} & \textbf{Aplic.} & \textbf{Situação observada} \\
		\midrule
		\endfirsthead
		\toprule
		\textbf{\#} & \textbf{Item} & \textbf{Aplic.} & \textbf{Situação observada} \\
		\midrule
		\endhead
		\bottomrule
		\endlastfoot
		1  & Requisito relacionado identificado & Sim & Todo artefato aponta para identificador \texttt{FR}, \texttt{NFR} ou \texttt{CA} \\
		\addlinespace
		2  & Critérios de aceitação verificados & Sim & 12 de 12 satisfeitos (Tabela~\ref{tab:ca-resultados}) \\
		\addlinespace
		3  & Estado do repositório inspecionado antes da alteração & Não & O diretório não era repositório válido no início da execução. Item não cumprido, registrado no \texttt{STATUS.md} \\
		\addlinespace
		4  & Alteração no menor escopo necessário & Sim & Cada tarefa restrita ao módulo afetado; nenhuma refatoração não relacionada \\
		\addlinespace
		5  & Compilação do alvo executada & Não & Não aplicável: artefato interpretado diretamente pelo navegador, sem etapa de compilação \\
		\addlinespace
		6  & Testes executados & Sim & 12 verificações por inspeção e execução observada \\
		\addlinespace
		7  & Análise estática executada quando aplicável & Sim & Busca textual de referências externas e de temporizadores fixos \\
		\addlinespace
		8  & Consumo de memória e tempo verificados quando aplicável & Não & Não medido. Registrado como dado necessário P-10 \\
		\addlinespace
		9  & Documentação atualizada & Sim & Documento de estado e documentação de uso produzidos \\
		\addlinespace
		10 & Rastreabilidade atualizada & Sim & Matriz do Apêndice~\ref{ap:rastreabilidade} \\
		\addlinespace
		11 & Arquivo de estado atualizado & Sim & Documento de estado com pendências e riscos residuais \\
		\addlinespace
		12 & Arquivo de transferência atualizado apenas se houver continuidade & Não & Não houve transferência a outro responsável na execução relatada \\
		\addlinespace
		13 & Riscos residuais registrados & Sim & Quatro riscos residuais, incluindo a pendência de validação em dispositivo físico \\
		\addlinespace
		14 & Nenhuma alteração não relacionada incorporada & Sim & Escopo preservado; nenhuma dependência adicionada \\
	\end{longtable}
}

\section{Leitura da lista}
\label{ap:checklist-leitura}

Dos catorze itens, nove foram cumpridos, quatro não se aplicam ao tipo de artefato
e um --- a inspeção do estado do repositório antes da alteração --- não foi
cumprido, por indisponibilidade do próprio mecanismo de controle de versão no
ambiente de execução.

A consequência prática do item 3 é relevante e merece registro. O portão G0 do
framework exige a leitura do estado do repositório antes de qualquer alteração,
precisamente para impedir que o executor trate uma versão lembrada de um arquivo
como se fosse a versão presente. Como o repositório não era válido no início da
execução, o portão não pôde ser cumprido, e a mitigação adotada foi a leitura
direta dos arquivos antes de cada alteração.

Isso expõe uma fragilidade do framework: um portão que depende de infraestrutura
externa não possui plano alternativo declarado. A correção proposta é que o portão
G0 defina explicitamente o procedimento substituto quando o sistema de controle de
versão não estiver disponível, em vez de simplesmente não ser cumprido.
